%!TEX root = ../main.tex
% Resumen en español (~250 palabras). Las palabras clave se definen
% en main.tex con \palabrasclave{...}

\begin{resumen}

La organización de un evento social o empresarial exige contratar por
separado servicios como catering, fotografía, decoración o ambientación,
usando distintas plataformas, redes sociales y medios de contacto, sin
herramientas que permitan anticipar problemas ni comparar alternativas
antes de decidir. Esta situación afecta a quienes organizan eventos y
también a los salones y proveedores, que gestionan sus clientes y
reservas de forma manual o con herramientas independientes.

Este trabajo desarrolla una plataforma que reúne en un mismo entorno un
marketplace de servicios para eventos, herramientas de gestión para
salones y proveedores, un sistema de recomendación basado en reglas
ponderadas con motivo explicable y un Gemelo Digital que permite simular
cambios en la cantidad de invitados, la fecha o el presupuesto sin
alterar el evento real. El sistema se construyó sobre el prototipo de
gestión de un caso real, Rodríguez Catering, con una arquitectura
multi-tenant, un servidor en Laravel y un cliente en React. La propuesta
se validó con una encuesta a potenciales usuarios, tres entrevistas a
especialistas del sector y una suite de 410 pruebas automatizadas, y se
evaluó su viabilidad económica en tres escenarios.

Los resultados muestran que es técnicamente viable integrar las tres
capas; que el 91,1\,\% de las personas encuestadas usaría o
probablemente usaría una plataforma de este tipo y que valoraron útiles la
recomendación (93,3\,\%) y la simulación (97,8\,\%); y que el modelo de
suscripción más comisión resulta viable en los tres escenarios, con valor
actual neto positivo entre USD~33.700 y USD~133.300. La mejora de
eficiencia frente al uso de herramientas fragmentadas no se midió y
queda como línea de trabajo futuro. El aporte consiste en un diseño
verificable que integra en una sola plataforma funciones que las
alternativas relevadas ofrecen por separado.

\end{resumen}
